\documentclass[10pt]{article}

\usepackage[utf8]{inputenc}

\usepackage[T1]{fontenc}

\usepackage{amsmath}

\usepackage{amsfonts}

\usepackage{amssymb}

\usepackage[version=4]{mhchem}

\usepackage{stmaryrd}



\begin{document}

Асимптотика: при фиксированном $v|\arg z|<3 \pi / 4$ и $|z| \rightarrow \infty$



\[

\begin{aligned}

D_{v}(z)=z^{v} e^{-z^{2} / 4} & {\left[\sum_{k=0}^{N} \frac{(-v / 2)_{k}(1 / 2-v / 2)_{k}}{k !}\left(\frac{z^{2}}{-2}\right)^{-k}+\right.} \\

& +O\left(|z|^{-2 N-2}\right),

\end{aligned}

\]



пріг ограниченном $|z|,|\arg (-v)| \leqslant \pi / 2$ и $|v| \rightarrow \infty$



\[

\begin{aligned}

D_{v}(z)=\frac{1}{\sqrt{2}} & \exp \left[\frac{v}{2} \ln (-v)-\frac{v}{2}-\sqrt{-v} z\right] \times \\

& \times\left[1+O\left(\frac{1}{\sqrt{|v|}}\right)\right] .

\end{aligned}

\]



П. ц. Ф. связана с др. функциями следующими соотношениями. С многочленами Эрмита:



\[

D_{n}(z)=2-n / 2 e^{-z^{2} / 4} H_{n}(z / \sqrt{2}), n=0,1,2, \ldots ;

\]



с интегралом вероятности:



\[

\begin{gathered}

D_{-n-1}(z)=\frac{(-1)^{n} \sqrt{2}}{n !} e^{-z^{2} / 4} \frac{d^{n}}{d z^{n}}\left(e^{z^{2} / 2} \operatorname{erf} \frac{z}{\sqrt{z}}\right), \\

n=0,1,2, \ldots ;

\end{gathered}

\]



с функциями Бесселя:



\[

D_{-1 / 2}(z)=\sqrt{\frac{\pi z}{2}} K_{1 / 4}\left(\frac{z^{2}}{4}\right) \text {. }

\]



Лuт.: [1] Б е й т м е н Г., Э цендентные Функции. Функции Бесселя, Функции параболического цилиндра, ортогональные многочлены, пер. с англ. 2 изд., М., 1974; [2] М и л л е p Д ж.-Ч.-П., Таблицы функций Вебера (Функций параболического цилиндра), пер. с англ.,



М., 1968 . Ю. А. Брычков, А. П. Прудников.



ПАРАБОЛОИД - незамкнутая нецентральная noверхность второго порядка. Канонич. уравнения П.:



\[

\frac{x^{2}}{p}+\frac{y^{2}}{q}=2 z, p, q>0

\]



\begin{itemize}

  \item эллиптический параболоид (при $p=q$ называется П. вращения) и

\end{itemize}



\[

\frac{x^{2}}{p}-\frac{y^{2}}{q}=2 z, p, q>0,

\]



\begin{itemize}

  \item гиперболический параболоид.

\end{itemize}



А. Б. Иванов.



ПАРАБОЛОИДАЛЬНЫЕ КООРДИНАТЫ - числа $u, v$ и $w$, связанные с прямоугольными координатами $x$, $y$ и $z$ формулами:



\[

x=2 u w \cos v, y=2 u w \sin v, z=u^{2}-w^{2},

\]



где $0 \leqslant u<\infty, 0 \leqslant v<2 \pi, 0 \leqslant w<\infty$. Координатные поверхности: две системы параболоидов вращения с иротивоположно направленными осями ( $u=$ const и $w=$ const) и полуплоскости $(v=$ const). Система П. к.ортогональная.



Коэффициенты Ламе:



\[

L_{u}=L_{w}=2 \sqrt{u^{2}+w^{2}}, L_{v}=2 u w .

\]



Элемент шлощади говерхности:



$d \sigma=4 \sqrt{\left(u^{2}+w^{2}\right) u^{2} w^{2}\left(d u^{2}+d w^{2}\right) d v^{2}+\left(u^{2}+w^{2}\right)(d u d w)^{2}}$.



Элемент объема:



\[

d V=8\left(u^{2}+w^{2}\right) \iota i w d u d v d w .

\]



Векторные дифференциальные операции:



\[

\begin{aligned}

\operatorname{grad}_{u} \varphi= & \frac{1}{2 \sqrt{u^{2}+w^{2}}} \frac{\partial \varphi}{\partial u}, \operatorname{grad}_{v} \varphi=\frac{1}{2 u w} \frac{\partial \varphi}{\partial v}, \\

& \operatorname{grad}_{w} \varphi=\frac{1}{2 \sqrt{u^{2}+w^{2}}} \frac{\partial \varphi}{\partial w},

\end{aligned}

\]



$\operatorname{div} \boldsymbol{c} \boldsymbol{\iota}=\frac{1}{2 u w \sqrt{\left(u^{2}+w^{2}\right)^{3}}}\left[w a_{u}\left(2 u^{2}+w^{2}\right)+u a_{w}\left(u^{2}+2 w^{2}\right)\right]+$ $+\frac{1}{2 \sqrt{u^{2}+w^{2}}}\left(\frac{\partial a_{u}}{\partial u}+\frac{\partial a_{w}}{\partial w}\right)+\frac{1}{2 u w} \frac{\partial a_{v}}{\partial v} ;$



\[

\begin{aligned}

& \operatorname{rot}_{u} \boldsymbol{a}=\frac{1}{2 u w} \frac{\partial a_{w}}{\partial v}-\frac{1}{2 w \sqrt{u^{2}+w^{2}}}\left(a_{v}+w \frac{\partial a_{v}}{\partial w}\right), \\

& \operatorname{rot}_{v} \boldsymbol{a}=\frac{1}{2 \sqrt{\left(u^{2}+w^{2}\right)^{3}}}\left(w a_{u}-a_{w}\right)+ \\

& +\frac{1}{2 \sqrt{u^{2}+w^{2}}}\left(\frac{\partial a_{u}}{\partial w}-\frac{\partial a_{w}}{\partial u}\right) \\

& \operatorname{rot}_{w} \boldsymbol{a}=\frac{1}{2 w\left(u^{2}+w^{2}\right)}\left(a_{v}+u \frac{\partial a_{v}}{\partial u}\right)-\frac{1}{2 u v} \frac{\partial a_{u}}{\partial v} ; \\

& \Delta \varphi=\frac{1}{4\left(u^{2}+w^{2}\right)}\left[\frac{\partial^{2} u}{\partial u^{2}}+\frac{1}{u} \frac{\partial \varphi}{\partial u}+\left(\frac{1}{u^{2}}+\frac{1}{w^{2}}\right) \frac{\partial^{2} \varphi}{\partial v^{2}}+\right. \\

& \left.+\frac{\partial^{2} \varphi}{\partial w^{2}}+\frac{1}{w} \frac{\partial \varphi}{\partial w}\right] \text {. }

\end{aligned}

\]



Д. Д. Соколов.



ПАРАДОКС - то же, что антиномия.



ПАРАКОМПАКТНОЕ ПРОСТРАНСТВО - топологическое пространство, в любое открытое покрытие к-рого можно вписать локально конечное открытое покрытие. (Семейство $\gamma$ множеств, лежащих в топологич. пространстве $X$, наз. л о к а л ь н о к о н е ч ны м в $X$, если у каждой точки $x \in X$ существует окрестность в $X$, пересекающаяся лишь с конечным множеством элементов семейства $\gamma$; семейство $\gamma$ множеств в пі и с ав о в семейство $\lambda$ множеств, если каждый элемент семейства $\gamma$ содержітся в нек-ром әлементе семейства $\lambda$. П а р а к о м ІІ а к т о м наз. паракомшактное хаусдорфово пространство. Класс паракомпактов весьма широк - он включает все метрич. пространства (теорема Стоуна) и все бикомпакты. Однако не каждое локально бикомпактное хаусдорфово Ігостранство паракомпактно.



Значение паракомпактности определяется отмеченной общностью этого понятия и рядом замечательных свойств паракомпактов. Прежце всего, каждое хаусдорфово П. П. нормально. Это позволяет строить на паракомпактах р а з б и е н и я е д и н и ц ы, подчиненные произвольному заданному открытому шокрытию $\gamma$. Так называются семейства действительных неотрицательных непрерывных функций ва пространстве, подчиненные следующим условиям: а) семейство носителей этих функций локально конечно и вписано в $\gamma$; б) в каждой точке пространства сумма значений всех тех функций семейства, к-рые отличны в ней от нуля (а таких функций конечное qисло), равна 1. Разбиения единицы являются основным средством построения погружений пространств в стандартные пространства. В частности, они используются для вложений многообразий в евклидовы пространства и при доказательстве теоремы о метризуемости каждого тихоновского пространства с $\sigma$-локально конечной базой. Кроме того, на разбиения $\mathbf{x}$ единицы в теории многообразий основаны методы, с помощью к-рых осуществляется единовременный синтез локальных шостроений, шроизводимых в ігеделах отдельных карт (в частности, тех или иных векторных и тензорных полей). Поэтому одним из требований исходных в теории многообразий является требование паракомшактности, не являюшееся лишним, т. к. существуют непаракомпактные связные хаусдорфовы многообразия.



В присутствии паракомпактности нек-рые локальные свойства пространства синтезируются и выполняются глобально. В частности, если паракомшакт локально метризуем, то он метризуем; если хаусдорфово пространство локально полно по Чеху и паракомпактно, то оно полно по Чеху. В размерности теории для паракомпактов удается получить ряд важных соотношений, не распространяющихся даже на нормальные пространства. Это не удивительно, т. к. одно из основных определений размерности - по ЈІебегу - связано с рассмотрением кратности открытых покрытий, что, несомненно,





\end{document}